\noindent\textbf{LLMs for Digital Twins and Simulation Generation.}
Recent work has begun to explore how LLMs can support digital-twin (DT) and simulation-related workflows.
In \cite{xia2024LLM_multi_agent_simulation}, the authors study LLMs for multi-agent DT simulations in a simplified 2D grid-world setting, and evaluate outputs via similarity between generated code and a reference implementation.
While this provides an initial validation, such similarity-based evaluation can under-specify functional correctness and may not transfer cleanly to simulator-grade DTs with richer physics, longer horizons, and more heterogeneous components.
Complementary efforts investigate LLM-driven decision-making in multi-agent simulations \cite{zhang2024llmDT_decision}, demonstrating the potential for automating decision processes in dynamic environments.
Additionally, \cite{mu2024rule} examines rule-based reward design for training LLMs in controlled settings, motivating the need for evaluation and reward signals that remain reliable as task complexity grows.

\noindent\textbf{Coding Benchmarks and Execution-Based Evaluation.}
A large body of work evaluates code-generation models via unit tests or execution-based correctness.
AlphaCode assesses generated programs against hidden test cases in competitive programming settings \cite{li2022competition}.
CodeT uses self-generated unit tests to score correctness of function implementations \cite{chen2022codet}.
Beyond similarity metrics such as BLEU \cite{Papineni2002BleuAM}, ROUGE \cite{lin-2004-rouge}, and CodeBLEU \cite{ren2020codebleu}, many benchmarks increasingly emphasize execution-based evaluation, where generated code is run and compared against test cases \cite{austin2021program, chen2021evaluating, huang-etal-2022-execution, Lai2022DS1000, hendrycksapps2021}.
Execution-based datasets such as ExeDS \cite{huang-etal-2022-execution} and ODEX \cite{wang2023executionbased} further highlight the importance of testing code in realistic execution environments.
More recently, multi-turn evaluation settings have been proposed to better reflect interactive programming workflows.
For example, MINT evaluates coding LLMs with tool usage and language feedback in multi-turn interactions \cite{wang2024mintevaluatingllmsmultiturn}, and \cite{zhang2024pybenchevaluatingllmagent} benchmarks LLM agents operating in a multi-agent setting with local file-system interactions.

\noindent\textbf{LLM-as-a-Judge.}
There is also growing interest in using LLMs as evaluators for generated outputs.
Prior work studies LLM judges for multi-turn conversations and chain-of-thought reasoning \cite{zheng2023LLMasJudge}.
However, several studies caution that LLM-based evaluation can be biased or unreliable under certain conditions, including when evaluating code or text generation \cite{abdelaziz2021AAAI_benchmark_code_understanding, wang2023llm_unfair_eval}.
Relatedly, \cite{yuan2024selfrewardinglanguagemodels} proposes a self-rewarding paradigm in which an LLM serves as both generator and evaluator, forming a closed-loop refinement process.
These findings motivate careful judge design, calibration, and validation, especially in high-stakes domains such as simulator-ready DT generation.

\noindent\textbf{Positioning of This Work.}
A comprehensive survey of deep learning methods for code generation and evaluation is provided in \cite{wu2022surveydeeplearningmodelscodes}.
In contrast to general-purpose code benchmarks, our focus is on simulator-grade DT construction, where artifacts may span hundreds to thousands of lines and correctness is multi-faceted (physics consistency, component completeness, parameter validity, and scenario logic).
We introduce a benchmark methodology that couples a curated DT-generation dataset with an interpretable, rubric-based judge model (J-LLM) tailored to DT code generation.
This design enables statistically grounded comparisons against similarity-based and execution-based evaluation, while providing diagnostic feedback beyond a single scalar pass/fail outcome.
